\documentclass[11pt,a4paper]{article}
\usepackage{fontspec}
\setmainfont{DejaVu Sans}[Scale=0.90]
\usepackage{amsmath,amssymb}
\usepackage[margin=2.2cm]{geometry}
\usepackage{booktabs,array,tabularx}
\usepackage{enumitem}
\usepackage{titlesec}
\usepackage{parskip}
\titlelabel{\thetitle.\quad}
\titleformat{\section}{\normalfont\large\bfseries}{\thesection.}{0.6em}{}
\titleformat{\subsection}{\normalfont\normalsize\bfseries}{\thesubsection}{0.6em}{}
\titlespacing*{\section}{0pt}{15pt}{5pt}
\titlespacing*{\subsection}{0pt}{11pt}{4pt}
\setlist[itemize]{leftmargin=1.4em,itemsep=1pt,topsep=3pt}
\begin{document}
\begin{center}
{\Large\bfseries Le gaz de phonons}\\[4pt]
{\normalsize Vibrations du réseau, quantification, fonction de distribution}\\[2pt]
{\small Notes de partie I --- 12 septembre 2026}
\end{center}

\vspace{4pt}


Ces notes établissent les objets sur lesquels repose la description cinétique du transport
thermique dans un cristal. Elles ne supposent aucune connaissance préalable de la physique du
solide. La chaîne conduisant de l'équation de transport à l'équation de Guyer--Krumhansl fait
l'objet des notes suivantes.

\section{Vibrations du réseau}

\subsection{La chaîne monoatomique}

Le modèle le plus simple d'un cristal est une chaîne d'atomes identiques de masse $M$, espacés de
$a$, reliés par des ressorts de raideur $C$. En notant $u_n$ le déplacement du $n$-ième atome par
rapport à sa position d'équilibre, la loi de Newton appliquée aux deux ressorts voisins donne
\[ M\,\ddot{u}_n = C\,(u_{n+1} - u_n) + C\,(u_{n-1} - u_n) = C\,(u_{n+1} + u_{n-1} - 2u_n). \]

Le système est linéaire et invariant par translation d'un pas de réseau. Ses solutions propres sont
donc des ondes planes,
\[ u_n(t) = A\,e^{i(kna - \omega t)} , \]
où $k$ est le nombre d'onde et $\omega$ la pulsation.

\subsection{Relation de dispersion}

En reportant l'onde plane dans l'équation du mouvement et en simplifiant le facteur commun, il
vient $-M\omega^2 = C\,(e^{ika} + e^{-ika} - 2) = 2C\,(\cos ka - 1)$, soit
\[ \boxed{\;\omega(k) = 2\sqrt{\frac{C}{M}}\;\left|\sin\frac{ka}{2}\right|\;} \]

Cette relation entre $\omega$ et $k$ s'appelle la \emph{relation de dispersion}. Elle contient
toute la dynamique du réseau, et elle diffère fondamentalement du cas d'un milieu continu, où
$\omega = vk$ serait linéaire pour tout $k$.

\subsection{Limite des grandes longueurs d'onde}

Pour $ka \ll 1$, le sinus se confond avec son argument et
\[ \omega \simeq a\sqrt{\frac{C}{M}}\;k = v\,k , \qquad v = a\sqrt{\frac{C}{M}} . \]

Le réseau se comporte alors comme un milieu continu, et $v$ est la vitesse du son. La structure
discrète du cristal ne se manifeste que lorsque la longueur d'onde devient comparable au pas du
réseau.

\subsection{Zone de Brillouin}

La relation de dispersion est périodique en $k$, de période $2\pi/a$. Deux valeurs de $k$ différant
de $2\pi/a$ décrivent exactement le même mouvement des atomes : entre deux atomes voisins,
l'information sur la phase est indiscernable.

Il suffit donc de considérer l'intervalle
\[ -\frac{\pi}{a} \;\le\; k \;\le\; \frac{\pi}{a} , \]
appelé \emph{première zone de Brillouin}. Toute valeur extérieure s'y ramène par soustraction d'un
multiple de $2\pi/a$.

Cette périodicité n'est pas une commodité de notation. Elle est à l'origine des processus umklapp,
qui gouvernent la résistance thermique intrinsèque d'un cristal.

\subsection{Vitesse de groupe}

L'énergie n'est pas transportée à la vitesse de phase $\omega/k$ mais à la vitesse de groupe
\[ v_g = \frac{d\omega}{dk} = a\sqrt{\frac{C}{M}}\;\cos\frac{ka}{2} . \]

Elle vaut $v$ au centre de la zone et s'annule au bord, en $k = \pi/a$. Les modes de bord de zone
ne transportent donc pas d'énergie : ils correspondent à une onde stationnaire où deux atomes
voisins vibrent en opposition de phase.

C'est la vitesse de groupe, non la vitesse de phase, qui intervient dans toute expression de flux.

\subsection{Chaîne diatomique, branches acoustique et optique}

Lorsque la cellule élémentaire contient deux atomes de masses différentes, la résolution donne deux
solutions pour chaque $k$.

La branche \emph{acoustique} part de $\omega = 0$ en $k = 0$ : les deux atomes de la cellule se
déplacent en phase, et le mode se réduit à une translation d'ensemble à grande longueur d'onde.
C'est elle qui porte le son, et l'essentiel du transport thermique.

La branche \emph{optique} conserve une pulsation finie en $k = 0$ : les deux atomes vibrent en
opposition de phase, le centre de masse restant fixe. Sa vitesse de groupe est faible, et sa
contribution au transport est modeste. Son nom vient de ce que, dans un cristal ionique, un tel
mouvement crée un dipôle oscillant susceptible de coupler au champ électromagnétique.

\subsection{Cas tridimensionnel}

Dans un cristal à trois dimensions comportant $p$ atomes par cellule élémentaire, chaque valeur du
vecteur d'onde $\mathbf{k}$ porte $3p$ modes : trois branches acoustiques --- une longitudinale,
deux transverses --- et $3p - 3$ branches optiques.

Le nitrure d'aluminium cristallise dans la structure wurtzite, dont la cellule élémentaire contient
quatre atomes. Il possède donc douze branches : trois acoustiques et neuf optiques.

Un mode est dès lors repéré par le couple $(\mathbf{k}, s)$, où $s$ désigne la branche. Toute somme
sur les modes porte sur les deux indices. Dans la suite, l'indice de branche est laissé implicite.

\section{Quantification}

\subsection{Le phonon}

Chaque mode de vibration est un oscillateur harmonique. La mécanique quantique impose que son
énergie ne prenne que des valeurs discrètes,
\[ E_n = \left(n + \tfrac{1}{2}\right)\hbar\omega , \qquad n = 0, 1, 2, \ldots \]

Le quantum d'énergie $\hbar\omega$ porte le nom de \emph{phonon}. L'entier $n$ est le nombre de
phonons présents dans le mode.

Un phonon n'est pas une particule au sens d'un constituant de la matière : c'est une excitation
collective du réseau, une quasi-particule. Le terme de gaz de phonons désigne l'ensemble de ces
excitations, traitées comme des porteurs d'énergie susceptibles de se propager et de se heurter.

\subsection{Quasi-moment}

À un phonon de vecteur d'onde $\mathbf{k}$ on associe la quantité $\hbar\mathbf{k}$, appelée
\emph{quasi-moment} ou moment cristallin.

Elle se comporte comme une quantité de mouvement dans les lois de conservation, mais elle n'en est
pas une : elle n'est définie qu'à un vecteur du réseau réciproque près, conséquence directe de la
périodicité établie plus haut. Cette réserve est ce qui rend possible les processus umklapp.

\subsection{Statistique de Bose--Einstein}

Les phonons sont des bosons : un nombre arbitraire d'entre eux peut occuper le même mode. Leur
nombre n'étant pas conservé --- un phonon peut être créé ou détruit par interaction avec le réseau
--- le potentiel chimique est nul. Le nombre moyen d'occupation à l'équilibre thermique est
\[ \boxed{\;n^0(\omega, T) = \frac{1}{\exp\!\left(\dfrac{\hbar\omega}{k_B T}\right) - 1}\;} \]

\subsection{Limites de température}

Pour $\hbar\omega \ll k_B T$, le développement de l'exponentielle donne
\[ n^0 \simeq \frac{k_B T}{\hbar\omega} , \qquad\text{donc}\qquad \hbar\omega\, n^0 \simeq k_B T . \]
Chaque mode porte alors une énergie $k_B T$ indépendante de sa pulsation : c'est l'équipartition
classique, et la capacité thermique sature.

Pour $\hbar\omega \gg k_B T$, l'occupation décroît exponentiellement, $n^0 \simeq e^{-\hbar\omega/k_BT}$.
Les modes de haute pulsation sont gelés et cessent de contribuer.

\subsection{Température de Debye}

La frontière entre ces deux comportements est fixée par la température de Debye,
\[ \Theta_D = \frac{\hbar\omega_{\max}}{k_B} , \]
où $\omega_{\max}$ est la pulsation la plus élevée du spectre. Au-dessus de $\Theta_D$ tous les
modes sont excités classiquement ; en dessous, une part croissante du spectre est gelée.

Pour le nitrure d'aluminium, $\Theta_D$ est de l'ordre de mille kelvins. À température ambiante, le
cristal est donc nettement en dessous de sa température de Debye, et une fraction appréciable des
modes de haute pulsation contribue peu.

\section{La fonction de distribution}

\subsection{Définition}

L'état du gaz de phonons est décrit par la fonction
\[ f(\mathbf{r}, \mathbf{k}, t) , \]
qui donne le nombre moyen de phonons du mode $\mathbf{k}$ présents au voisinage du point
$\mathbf{r}$ à l'instant $t$.

À l'équilibre thermique uniforme, elle se réduit à la distribution de Bose--Einstein,
$f = n^0(\omega_\mathbf{k}, T)$, indépendante de $\mathbf{r}$. Tout écart à cette forme signale un
déséquilibre, et c'est de cet écart que naît le transport.

\subsection{Validité de la description}

Parler simultanément de la position et du vecteur d'onde d'un phonon suppose de construire des
paquets d'ondes localisés, ce qui n'est légitime que si la longueur d'onde reste petite devant
l'échelle sur laquelle les grandeurs macroscopiques varient.

Cette condition est la limite de validité de toute la description cinétique. Elle cesse d'être
remplie lorsque l'épaisseur d'un film devient comparable au libre parcours moyen des phonons,
régime dit balistique, où le transport ne se décrit plus par une équation de diffusion.

\subsection{Distribution déplacée}

Une seconde forme d'équilibre joue un rôle central dans la suite. Si le gaz de phonons possède une
vitesse d'ensemble $\mathbf{u}$, la distribution qui annule les collisions conservant le
quasi-moment est
\[ f^{\,\mathrm{d}} = \frac{1}{\exp\!\left(\dfrac{\hbar\omega - \hbar\mathbf{k}\cdot\mathbf{u}}{k_B T}\right) - 1} . \]

Elle décrit un gaz en équilibre interne mais globalement en mouvement, exactement comme une
maxwellienne déplacée décrit un fluide en écoulement. Sa présence est ce qui donne au gaz de
phonons un caractère hydrodynamique.

\section{Les moments de la distribution}

Les grandeurs mesurables ne sont pas la fonction de distribution elle-même, mais ses intégrales sur
l'espace des vecteurs d'onde, appelées \emph{moments}.

\subsection{Densité d'énergie}

\[ u(\mathbf{r}, t) = \int \frac{d^3k}{(2\pi)^3}\; \hbar\omega_\mathbf{k}\; f(\mathbf{r}, \mathbf{k}, t) \]

Chaque mode apporte son énergie $\hbar\omega$ multipliée par le nombre de phonons présents. C'est le
moment d'ordre zéro.

\subsection{Flux de chaleur}

\[ \mathbf{q}(\mathbf{r}, t) = \int \frac{d^3k}{(2\pi)^3}\; \hbar\omega_\mathbf{k}\;
\mathbf{v}_g(\mathbf{k})\; f(\mathbf{r}, \mathbf{k}, t) \]

Chaque phonon transporte son énergie à sa vitesse de groupe. C'est le moment d'ordre un.

À l'équilibre, $f$ ne dépend que de $\omega$, qui est une fonction paire de $\mathbf{k}$, tandis que
$\mathbf{v}_g$ est impaire. L'intégrale s'annule : \textbf{un gaz de phonons à l'équilibre ne
transporte aucune chaleur}. Le flux mesure exactement l'écart à l'équilibre.

\subsection{Quasi-moment et capacité}

Deux autres moments interviennent. La densité de quasi-moment,
\[ \mathbf{P} = \int \frac{d^3k}{(2\pi)^3}\; \hbar\mathbf{k}\; f , \]
est la grandeur que les processus normaux conservent.

Pour une dispersion linéaire isotrope $\omega = vk$, la vitesse de groupe vaut
$\mathbf{v}_g = v\,\hat{\mathbf{k}}$ et l'on obtient
\[ \mathbf{q} = v^2\,\mathbf{P} . \]

\textbf{Flux de chaleur et quasi-moment sont alors proportionnels.} Cette identité est ce qui rend
la suite tractable : une collision qui conserve le quasi-moment conserve du même coup le flux de
chaleur, et ne contribue donc pas à la résistance thermique.

La capacité thermique volumique s'obtient enfin par dérivation de la densité d'énergie,
\[ C = \frac{\partial u}{\partial T} . \]

\subsection{Ce qui manque}

Les moments sont les grandeurs mesurables, mais aucune équation ne les gouverne directement. Les
obtenir exige une équation d'évolution pour $f$, puis un moyen de fermer le système qu'elle engendre.
C'est l'objet des notes suivantes.

\section{Ordres de grandeur}

\begin{center}
\begin{tabularx}{\textwidth}{@{}lXl@{}}
\toprule
\textbf{Grandeur} & \textbf{Expression ou valeur} & \textbf{Ordre} \\
\midrule
Pas de réseau & distance entre atomes voisins & $3\times10^{-10}$ m \\
\addlinespace
Vitesse du son, AlN & branche longitudinale & $\sim 1{,}1\times10^{4}$ m/s \\
& branches transverses & $\sim 6\times10^{3}$ m/s \\
\addlinespace
Pulsation maximale & $\omega_{\max} = k_B\Theta_D/\hbar$ & $\sim 10^{14}$ rad/s \\
\addlinespace
Température de Debye, AlN & & $\sim 10^{3}$ K \\
\addlinespace
Temps de relaxation à 300 K & entre collisions de phonons & $10^{-12}$ à $10^{-11}$ s \\
\addlinespace
Libre parcours moyen & $\Lambda = v\,\tau$ & $10^{-8}$ à $10^{-7}$ m \\
\addlinespace
Capacité volumique, AlN & $\rho c_p$ & $2{,}4\times10^{6}$ J\,m$^{-3}$K$^{-1}$ \\
\addlinespace
Conductivité, AlN massif & & $150$ à $250$ W\,m$^{-1}$K$^{-1}$ \\
\bottomrule
\end{tabularx}
\end{center}

Le rapprochement du libre parcours moyen, de l'ordre de la dizaine de nanomètres, et de l'épaisseur
d'un film mince, de l'ordre de quelques centaines de nanomètres, situe le problème : les deux
échelles ne sont séparées que d'un facteur dix à cent. La description purement diffusive n'y est
donc pas acquise.
\end{document}